%% kapa: exam=1 round=25 year=2014 period=1 form=A subject=economics subject_ko=경제원론 seq=37 number=77
%% source: Q-Net 감정평가사 기출문제 (cst00309) artlSeq=5101207 file=extracted/1008/5101207/2014년 제25회 감정평가사 1차시험 기출문제/1교시(A형).hwp
%% conversion: kordoc 4.24.51 -> tools/gen1.py (auto); review=proofread-sonnet; flags=none
\begin{question}{77}
A재와 B재만 생산하는 국가의 재화 생산량과 가격이 다음과 같다. \par
{\small\centering
\begin{tabular}{|c|c|c|c|c|}
\hline
\makecell{연도} & \makecell{A가격} & \makecell{A생산량} & \makecell{B가격} & \makecell{B생산량} \\
\hline
\makecell{2012} & \makecell{( ㄱ )} & \makecell{100} & \makecell{12} & \makecell{50} \\
\hline
\makecell{2013} & \makecell{10} & \makecell{80} & \makecell{10} & \makecell{( ㄴ )} \\
\hline
\end{tabular}\par}

2012년도를 기준연도로 할 때 2013년의 실질GDP 성장률이 25 \%이고 GDP 디플레이터가 90이라면 ( ㄱ )과 ( ㄴ )에 알맞은 것은? \par
\begin{choices}
  \item ㄱ: 8, ㄴ: 60
  \item ㄱ: 8, ㄴ: 80
  \item ㄱ: 10, ㄴ: 80
  \item ㄱ: 10, ㄴ: 100
  \item ㄱ: 12, ㄴ: 100
\end{choices}
\end{question}
\answer{4}
